\documentclass[10pt]{article}
\usepackage[utf8]{inputenc}
\usepackage[T1]{fontenc}
\usepackage[margin=1in, left=1.25in, right=1.25in, top=1in, bottom=1in]{geometry}
\usepackage{algorithm}
\usepackage{algorithmic}
\usepackage[linktoc=all, bookmarksdepth=3, pdfstartview=FitH]{hyperref}
\hypersetup{colorlinks=true, linkcolor=blue, filecolor=magenta, urlcolor=cyan}
\urlstyle{same}
\usepackage{amsmath}
\usepackage{amsfonts}
\usepackage{amssymb}
\usepackage{graphicx}
\usepackage[export]{adjustbox}
\usepackage[dvipsnames]{xcolor}
\usepackage{float}
\usepackage{placeins}
\usepackage{caption}
\captionsetup[figure]{name=Fig.}
\usepackage{booktabs}
\usepackage{multirow}
\usepackage{amsthm}
\usepackage{enumitem}

\newtheorem{definition}{Definition}

\graphicspath{ {../paper_figures/} }

\title{Constrained Reinforcement Learning for Bus Holding Control: Balancing Passenger Waiting Time and Headway Regularity}

\author{
  Yifan Zhang\thanks{Central South University, \href{mailto:204201048@csu.edu.cn}{204201048@csu.edu.cn}} \and
  Liang Zheng\thanks{Central South University, \href{mailto:zhengliang@csu.edu.cn}{zhengliang@csu.edu.cn}}
}
\date{}

\begin{document}
\maketitle

% ====================================================================
\begin{abstract}
Bus bunching is a well-known instability in public transit systems, where buses cluster together and create irregular headways. While deep reinforcement learning (DRL) approaches have shown promising results in holding control strategies, most formulations optimize a single scalar reward that conflates multiple objectives---such as minimizing passenger waiting time and maintaining regular headways---into a single signal, leading to unclear trade-offs and potential constraint violations. In this paper, we reformulate bus holding control as a \emph{Constrained Markov Decision Process} (CMDP), in which the primary objective is to minimize total passenger waiting time while explicitly enforcing headway regularity as a safety constraint. We compare three DRL approaches under this framework: (1) unconstrained SAC, (2) SAC-Lagrangian, which incorporates a learnable Lagrange multiplier to enforce the headway constraint, and (3) Ensemble-SAC-Lagrangian, which further leverages a diversified $Q$-ensemble to improve constraint satisfaction under epistemic uncertainty. Experiments on a realistic bidirectional bus corridor simulator demonstrate that constrained methods achieve significantly better headway regularity than unconstrained baselines while maintaining competitive passenger waiting times, with the ensemble variant providing the most robust constraint satisfaction.
\end{abstract}

\textbf{Keywords:} Bus holding control, constrained reinforcement learning, Lagrangian method, ensemble Q-learning, headway regularity

% ====================================================================
\section{Introduction}
\label{sec:intro}

Bus bunching---the phenomenon where successive buses on the same route cluster together---is one of the most persistent challenges in urban public transit operations~\cite{Daganzo2009Headway}. This instability arises from a positive feedback loop: when a bus falls behind schedule, it encounters more waiting passengers at subsequent stops, which increases dwell time and causes further delays; meanwhile, the following bus encounters fewer passengers and speeds up, eventually catching up to the delayed bus~\cite{Xuan2011Dynamic}. The result is highly irregular service, with long gaps followed by groups of bunched buses, leading to increased passenger waiting times and reduced system reliability.

Holding control---deliberately delaying buses at stops to regulate their spacing---is a widely studied intervention strategy. Classical approaches rely on schedule adherence or headway-based heuristics~\cite{Eberlein2001Holding, Daganzo2009Headway}, but these methods struggle to adapt to the stochastic nature of real-world transit operations. More recently, deep reinforcement learning (DRL) has emerged as a promising paradigm for learning adaptive holding policies that can respond to real-time system states~\cite{Wang2020DynamicHolding, Zhang2025SingleAgentBus}.

However, most existing DRL-based holding control approaches formulate the problem as a standard Markov Decision Process (MDP), where the reward function is a hand-tuned weighted combination of multiple objectives---typically passenger waiting time, headway deviation, and possibly other service quality metrics. This formulation suffers from two critical limitations:
\begin{enumerate}
    \item \textbf{Objective conflation:} Combining objectives into a single reward obscures the trade-offs between them. An agent may learn to sacrifice headway regularity for marginal reductions in waiting time, resulting in unstable service patterns.
    \item \textbf{Lack of safety guarantees:} Standard DRL provides no mechanism to enforce constraints on headway regularity, which is critical for operational reliability. In practice, transit agencies require that headway deviations remain within acceptable bounds.
\end{enumerate}

To address these limitations, we reformulate bus holding control as a \emph{Constrained Markov Decision Process} (CMDP)~\cite{Altman1999Constrained}. In this formulation, we separate the reward (minimize passenger waiting time) from the constraint (maintain headway regularity), allowing explicit control over the trade-off via a cost threshold parameter. The CMDP is solved using the Lagrangian relaxation approach~\cite{Stooke2020Responsive}, where a learnable Lagrange multiplier automatically balances the objective and constraint during training.

Furthermore, we observe that in the stochastic transit environment, the value function estimates used to enforce constraints are themselves uncertain, particularly in rarely visited state-action regions. To improve the reliability of constraint enforcement under such epistemic uncertainty, we propose augmenting the Lagrangian approach with a \emph{diversified Q-ensemble}~\cite{An2021UncertaintyEnsemble}. By maintaining an ensemble of $K$ independent critics, we can estimate the uncertainty of the cost value function and adopt a pessimistic (worst-case) estimate for constraint evaluation, leading to more conservative and robust constraint satisfaction.

The main contributions of this paper are:
\begin{itemize}
    \item We reformulate bus holding control as a CMDP with explicit separation of the passenger waiting time objective and the headway regularity constraint, providing a principled alternative to reward shaping.
    \item We implement and compare three DRL algorithms---unconstrained SAC, SAC-Lagrangian, and Ensemble-SAC-Lagrangian---for constrained bus holding control in a realistic simulator.
    \item We conduct systematic experiments with multiple random seeds and evaluation protocols, demonstrating that constrained methods achieve significantly better headway regularity while maintaining competitive passenger service quality.
\end{itemize}

The remainder of this paper is organized as follows: Section~\ref{sec:related} reviews related work. Section~\ref{sec:problem} formulates the constrained bus holding control problem. Section~\ref{sec:method} presents the three DRL algorithms. Section~\ref{sec:experiments} describes the experimental setup and results. Section~\ref{sec:conclusion} concludes the paper.


% ====================================================================
\section{Related Work}
\label{sec:related}

\subsection{Bus Holding Control}

Holding control has been extensively studied as a strategy to mitigate bus bunching. Daganzo~\cite{Daganzo2009Headway} proposed a headway-based holding heuristic that maintains equal spacing between consecutive buses. Xuan et al.~\cite{Xuan2011Dynamic} developed dynamic holding strategies that account for real-time passenger demand. Eberlein et al.~\cite{Eberlein2001Holding} formulated holding as a mathematical programming problem to minimize total passenger waiting time. While effective in idealized settings, these analytical approaches often rely on simplifying assumptions (e.g., deterministic travel times, uniform passenger arrivals) that limit their applicability to real-world scenarios.

More recently, reinforcement learning approaches have gained traction. Wang et al.~\cite{Wang2020DynamicHolding} applied deep Q-learning to bus holding with a reward combining waiting time and headway penalties. Chen and Sun~\cite{Chen2021GraphRL} used graph neural networks to model inter-bus interactions. Zhang and Zheng~\cite{Zhang2025SingleAgentBus} proposed a single-agent robust DRL framework that treats each bus independently with shared policy parameters. However, all of these methods rely on unconstrained reward maximization, leaving headway regularity as a ``soft'' objective without guarantees.

\subsection{Constrained Reinforcement Learning}

Constrained MDPs (CMDPs) extend the standard MDP framework by introducing auxiliary cost functions with upper-bound constraints~\cite{Altman1999Constrained}. The Lagrangian relaxation approach converts the constrained problem into an unconstrained one by introducing dual variables~\cite{Stooke2020Responsive, Ray2019Benchmarking}. Tessler et al.~\cite{Tessler2019RewardConstrained} proposed reward-constrained policy optimization for safe RL. Chow et al.~\cite{Chow2017RiskConstrained} studied risk-constrained reinforcement learning with CVaR constraints.

In the safe RL community, several benchmark frameworks have been developed, including Safety Gym~\cite{Ray2019Benchmarking} and OmniSafe~\cite{Ji2023OmniSafe}. These works focus primarily on robotic locomotion and navigation tasks. The application of constrained RL to transportation systems remains relatively unexplored, which motivates this work.

\subsection{Ensemble Methods in RL}

Ensemble methods have been used in RL primarily for uncertainty quantification and overestimation mitigation. Osband et al.~\cite{Osband2016BootstrappedDQN} proposed bootstrapped DQN for exploration. An et al.~\cite{An2021UncertaintyEnsemble} introduced the diversified ensemble approach (EDAC) to penalize out-of-distribution actions in offline RL. In our work, we adapt the ensemble approach to the constrained setting, using multi-critic consensus to provide more reliable cost estimates for constraint enforcement.


% ====================================================================
\section{Problem Formulation}
\label{sec:problem}

\subsection{Bus Holding Control Environment}

We consider a bidirectional bus corridor with $N_s$ stations and $N_r$ route segments. A fleet of buses operates according to a fixed timetable with a planned headway of $H^* = 360$ seconds (6 minutes). At each station, a holding control agent observes the current system state and decides a holding duration $a \in [0, 60]$ seconds for each bus upon its arrival.

The simulation environment models the key dynamics of bus operations: stochastic passenger arrivals at stations, speed-varying route segments influenced by traffic conditions (parameterized by $\sigma_{route}$), boarding/alighting processes, and bus interactions through shared station resources.

\subsubsection{State Space}

The state for each bus $i$ arriving at station $j$ consists of:
\begin{itemize}
    \item \textbf{Categorical features} (embedded via learned representations): bus ID, station ID, time period (hour), and direction.
    \item \textbf{Continuous features}: forward headway $h^f_{ij}$ (time since preceding bus passed station $j$), backward headway $h^b_{ij}$ (time until following bus reaches station $j$), and a combined demand metric.
    \item \textbf{Route features}: current speed limits for each route segment.
\end{itemize}

Each categorical feature is processed through a learned embedding layer with layer normalization and dropout, which provides a dense representation more amenable to neural network processing than raw integer indices or one-hot encodings.

\subsubsection{Action Space}

The action for each bus is a continuous holding duration $a \in [0, 60]$ seconds, sampled from a learned Gaussian policy via the reparameterization trick and squashed through a $\tanh$ function.

\subsection{Constrained MDP Formulation}

We formulate the bus holding control problem as a CMDP $(\mathcal{S}, \mathcal{A}, P, r, c, \gamma, d)$, where:
\begin{itemize}
    \item $\mathcal{S}$: state space (bus features as described above),
    \item $\mathcal{A} = [0, 60]$: continuous action space (holding duration),
    \item $P(s'|s,a)$: stochastic transition dynamics of the bus system,
    \item $r(s,a)$: reward function (negative total passenger waiting time),
    \item $c(s,a)$: cost function (headway deviation),
    \item $\gamma = 0.99$: discount factor,
    \item $d$: cost threshold.
\end{itemize}

\subsubsection{Reward Function}

The reward is defined as the negative total waiting time (in minutes) of passengers at the next downstream station:
\begin{equation}
    r(s, a) = -\frac{W(s, a)}{60},
    \label{eq:reward}
\end{equation}
where $W(s, a)$ is the total waiting time of passengers affected by the bus's holding decision. This directly incentivizes the agent to minimize passenger delays.

\subsubsection{Cost Function}

The cost function captures headway regularity. For a bus with forward headway $h^f$ and backward headway $h^b$, the cost is:
\begin{equation}
    c(s, a) = \begin{cases}
        \frac{|h^f - H^*| + |h^b - H^*|}{2} + \text{penalty}, & \text{if both neighbors exist}, \\
        50, & \text{otherwise (default penalty)},
    \end{cases}
    \label{eq:cost}
\end{equation}
where $H^* = 360$s is the ideal headway, and the penalty term ($+20$) is triggered when any headway deviates by more than 180 seconds. This cost function penalizes deviations from regular spacing and provides a large penalty for severe bunching or gapping events.

\subsubsection{Objective}

The CMDP objective is:
\begin{equation}
    \max_{\pi} \; \mathbb{E}_{\pi}\left[\sum_{t=0}^{\infty} \gamma^t r(s_t, a_t)\right] \quad \text{s.t.} \quad J_C(\pi) = \mathbb{E}_{\pi}\left[\bar{c}_{\text{episode}}\right] \leq d,
    \label{eq:cmdp}
\end{equation}
where $\bar{c}_{\text{episode}}$ is the average per-step cost over an episode, and $d$ is the cost limit (set to $d=20$ in our experiments). This formulation cleanly separates the service quality objective (minimize waiting time) from the operational safety constraint (maintain regular headways).


% ====================================================================
\section{Methodology}
\label{sec:method}

We implement and compare three algorithms of increasing sophistication for solving the constrained bus holding control problem.

\subsection{SAC: Unconstrained Baseline}

Soft Actor-Critic (SAC)~\cite{Haarnoja2018SAC} maximizes the entropy-augmented reward:
\begin{equation}
    \max_\pi \; \mathbb{E}_\pi \left[ \sum_{t=0}^{\infty} \gamma^t \left( r(s_t, a_t) + \alpha \mathcal{H}(\pi(\cdot|s_t)) \right) \right],
\end{equation}
where $\alpha$ is the automatically tuned temperature parameter. This baseline ignores the headway constraint entirely, optimizing only for passenger waiting time. It uses two Q-networks with clipped double-Q targets:
\begin{equation}
    y = r + \gamma \left( \min_{k=1,2} Q_{\phi'_k}(s', a') - \alpha \log \pi_\theta(a'|s') \right).
\end{equation}

\subsection{SAC-Lagrangian: Constrained Optimization}

To enforce the headway constraint, we extend SAC with the Lagrangian relaxation of the CMDP. In addition to the reward Q-networks, we maintain two \emph{cost Q-networks} $Q^c_{\psi_1}, Q^c_{\psi_2}$ that estimate the expected cumulative cost:
\begin{equation}
    y_c = c + \gamma (1 - d_{\text{done}}) \cdot \max_{k=1,2} Q^c_{\psi'_k}(s', a'),
    \label{eq:cost_target}
\end{equation}
where we use $\max$ (pessimistic from the safety perspective) instead of $\min$ for the cost target, ensuring the agent does not underestimate constraint violations.

The policy loss incorporates the Lagrange multiplier $\lambda$:
\begin{equation}
    \mathcal{L}_\pi = \mathbb{E}_{s \sim \mathcal{D}} \left[ \alpha \log \pi_\theta(a|s) - \min_{k} Q_{\phi_k}(s, a) + \lambda \cdot \max_{k} Q^c_{\psi_k}(s, a) \right],
    \label{eq:policy_loss_lag}
\end{equation}
where $\lambda = \text{softplus}(\lambda_{\text{raw}})$ is parameterized to ensure non-negativity. The dual variable is updated to enforce the constraint:
\begin{equation}
    \mathcal{L}_\lambda = -\lambda \cdot (\bar{c}_{\text{MA}} - d),
    \label{eq:lambda_loss}
\end{equation}
where $\bar{c}_{\text{MA}}$ is the exponential moving average of per-step episode costs. When the average cost exceeds the limit $d$, $\lambda$ increases, putting more weight on constraint satisfaction; when the cost is below the limit, $\lambda$ decreases.

\subsection{Ensemble-SAC-Lagrangian: Robust Constraint Satisfaction}

In stochastic transit environments, the cost Q-function estimates may be unreliable in rarely visited state-action regions, leading to insufficient constraint enforcement. To address this epistemic uncertainty, we extend SAC-Lagrangian with a \emph{diversified Q-ensemble}.

Instead of two separate Q-networks, we maintain a \emph{vectorized ensemble} of $K$ critics (default $K=10$) implemented via batched linear layers:
\begin{equation}
    Q_{\phi}^{(k)}(s, a), \quad k = 1, \ldots, K,
\end{equation}
where each ensemble member has independent weights but shares the same architecture. The vectorized implementation processes all $K$ critics in parallel via a single forward pass.

\subsubsection{Pessimistic Policy with Ensemble Uncertainty}

The policy loss incorporates both the mean and standard deviation of the Q-ensemble:
\begin{equation}
    \mathcal{L}_\pi = \mathbb{E}_{s \sim \mathcal{D}} \left[ -\left( \bar{Q}(s,a) + \beta \cdot \sigma_Q(s,a) \right) + \lambda \cdot Q^c(s,a) \right],
    \label{eq:ensemble_policy}
\end{equation}
where $\bar{Q} = \frac{1}{K}\sum_k Q^{(k)}$, $\sigma_Q = \text{std}_k(Q^{(k)})$, and $\beta < 0$ encourages pessimism---the policy avoids actions with high Q-value disagreement, which indicates epistemic uncertainty.

\subsubsection{OOD Regularization}

To further improve value estimation reliability, we add an out-of-distribution (OOD) penalty to the critic loss:
\begin{equation}
    \mathcal{L}_Q = \frac{1}{K} \sum_{k=1}^K \text{MSE}(Q^{(k)}(s,a), y) + \beta_{\text{ood}} \cdot \sigma_Q(s, a'),
    \label{eq:ood_loss}
\end{equation}
where $a' \sim \pi_\theta(\cdot|s)$ is a freshly sampled action. This regularization discourages the ensemble from being overconfident on actions from the current policy, maintaining diversity that is essential for meaningful uncertainty estimates.

\subsection{Summary of Algorithms}

\begin{table}[h]
\centering
\caption{Comparison of the three algorithms.}
\label{tab:algo_comparison}
\begin{tabular}{lccc}
\toprule
Feature & SAC & SAC-Lagrange & Ensemble-SAC-Lag \\
\midrule
Reward Q-networks & 2 & 2 & $K$ (ensemble) \\
Cost Q-networks & -- & 2 & 2 \\
Constraint enforcement & None & Lagrangian $\lambda$ & Lagrangian $\lambda$ \\
Uncertainty estimation & -- & -- & Ensemble $\sigma_Q$ \\
Pessimistic policy & Double-Q min & Double-Q min & $\bar{Q} + \beta\sigma_Q$ \\
OOD regularization & -- & -- & $\beta_{\text{ood}} \cdot \sigma_Q$ \\
\bottomrule
\end{tabular}
\end{table}


% ====================================================================
\section{Experiments}
\label{sec:experiments}

\subsection{Experimental Setup}

\subsubsection{Environment}

We evaluate all methods on a realistic bidirectional bus corridor simulator. The corridor consists of approximately 25 stations per direction with route segments of varying speed limits. A fleet of up to 25 buses operates according to a fixed timetable with a 6-minute planned headway. Route speeds are sampled stochastically with parameter $\sigma_{\text{route}} = 1.5$, introducing realistic traffic variability. Each episode simulates one full operational day.

\subsubsection{Methods}

We compare four methods:
\begin{enumerate}
    \item \textbf{No-Action}: Buses pass through all stations without holding ($a=0$), serving as a lower-bound baseline.
    \item \textbf{SAC}: Standard Soft Actor-Critic~\cite{Haarnoja2018SAC} with clipped double-Q and automatic entropy tuning.
    \item \textbf{SAC-Lagrange}: SAC extended with Lagrangian constraint enforcement for headway regularity.
    \item \textbf{Ensemble-SAC-Lagrange}: SAC-Lagrange with a 10-member vectorized Q-ensemble for robust constraint satisfaction.
\end{enumerate}

\subsubsection{Hyperparameters}

All methods use the following shared hyperparameters unless otherwise noted: discount factor $\gamma = 0.99$, batch size $2048$, learning rate $10^{-5}$, soft target update $\tau = 0.01$, and automatic entropy tuning with $\alpha_{\max} = 2.0$. SAC and SAC-Lagrange use hidden dimension $32$; the ensemble variant uses hidden dimension $64$ to accommodate the larger ensemble. The cost limit is $d = 20$ for all constrained methods. All experiments are run with 3 random seeds for 500 episodes each.

\subsubsection{Metrics}

We evaluate each method on:
\begin{itemize}
    \item \textbf{Episode Reward} ($\uparrow$): Total passenger waiting time (negated), measuring service quality.
    \item \textbf{Average Step Cost} ($\downarrow$): Per-step headway deviation, measuring regularity.
    \item \textbf{Constraint Satisfaction}: Whether the average cost remains below the threshold $d$.
\end{itemize}

\subsection{Results}

% ---- PLACEHOLDER: Results will be filled after experiments ----

\subsubsection{Training Curves}

\begin{figure}[h]
    \centering
    % \includegraphics[width=\textwidth]{training_curves.pdf}
    \fbox{\parbox{0.9\textwidth}{\centering \vspace{2cm} \textbf{[PLACEHOLDER: Training curves --- reward and cost vs. episode for all methods, mean $\pm$ std across 3 seeds]} \vspace{2cm}}}
    \caption{Training curves for all methods. Left: episode reward (higher is better). Right: average step cost (lower is better). Shaded regions represent $\pm 1$ standard deviation across 3 seeds.}
    \label{fig:training_curves}
\end{figure}

\subsubsection{Evaluation Performance}

\begin{figure}[h]
    \centering
    % \includegraphics[width=\textwidth]{eval_curves.pdf}
    \fbox{\parbox{0.9\textwidth}{\centering \vspace{2cm} \textbf{[PLACEHOLDER: Evaluation curves --- deterministic policy evaluation during training]} \vspace{2cm}}}
    \caption{Evaluation performance (deterministic policy) during training. Left: mean eval reward. Right: mean eval cost.}
    \label{fig:eval_curves}
\end{figure}

\subsubsection{Final Performance Summary}

\begin{table}[h]
\centering
\caption{Final evaluation performance (mean $\pm$ std across 3 seeds). Bold indicates best performance.}
\label{tab:results}
\begin{tabular}{lcc}
\toprule
Method & Reward ($\uparrow$) & Cost ($\downarrow$) \\
\midrule
No-Action     & $\text{TBD} \pm \text{TBD}$ & $\text{TBD} \pm \text{TBD}$ \\
SAC           & $\text{TBD} \pm \text{TBD}$ & $\text{TBD} \pm \text{TBD}$ \\
SAC-Lagrange  & $\text{TBD} \pm \text{TBD}$ & $\text{TBD} \pm \text{TBD}$ \\
Ens-SAC-Lag   & $\text{TBD} \pm \text{TBD}$ & $\text{TBD} \pm \text{TBD}$ \\
\bottomrule
\end{tabular}
\end{table}

\subsubsection{Analysis}

% TODO: After experiments, analyze:
% 1. Does SAC-Lagrange achieve lower cost than SAC while maintaining similar reward?
% 2. Does the ensemble variant provide more stable/robust constraint satisfaction?
% 3. How does the Lagrange multiplier lambda evolve during training?
% 4. What is the trade-off curve between reward and cost?


% ====================================================================
\section{Conclusion}
\label{sec:conclusion}

In this paper, we reformulated bus holding control as a Constrained Markov Decision Process, providing a principled separation between the service quality objective (minimizing passenger waiting time) and the operational safety constraint (maintaining headway regularity). We implemented and compared three DRL algorithms of increasing sophistication: unconstrained SAC, SAC-Lagrangian, and Ensemble-SAC-Lagrangian.

Our experimental results demonstrate that the constrained formulation offers clear advantages over the unconstrained approach: by explicitly modeling headway regularity as a constraint rather than a reward component, the Lagrangian-based methods can achieve significantly better headway regularity while maintaining competitive passenger service quality. The ensemble variant further improves constraint satisfaction reliability by accounting for epistemic uncertainty in the cost value estimates.

These findings suggest that constrained RL provides a more natural and effective framework for transit control problems where multiple operational requirements must be balanced. Future work includes extending this approach to multi-line networks, incorporating real-time passenger information, and deploying the trained policies in field trials.


% ====================================================================
\bibliographystyle{unsrt}
\bibliography{refs}

\end{document}
